\subsection{Baumdiagramme}\index{Baumdiagramm!Stochastik}

\textbf{Beispiel: Glücksrad}\index{Glücksrad}

%%\subsubsection{Einstiegsbeispiel: Glücksrad}
\bbwCenterGraphic{4cm}{geso/stoch/img/gluecksrad75deg}
An einem Glücksrad wird gedreht. $75\degre$ erzielen sofort einen Gewinn. Die anderen $285\degre$ geben einen Verlust an. Zum Glück darf ich am Rad drei mal drehen. Das Spiel endet entweder beim ersten Gewinn oder aber spätestens nach dreimaligem Drehen.

Wie groß ist die Wahrscheinlichkeit, bei diesem Spiel zu gewinnen?

Zeichnen Sie dazu das Baumdiagramm, das bei jedem ``Sieg'' sofort endet und
maximal drei Stufen tief geht.

\TNTeop{Einseitiger entarteter Baum mit Sieg = 75/360 und Verlust =
  285/360. 
  \bbwCenterGraphic{7cm}{geso/stoch/img/Rad75Baum.png}

  (Gegenereignis: Verlust bei $\approx 49.6\%$; somit Gewinn beim Gegenereignis.)

  $$P(\text{Gewinn}) = 1 - P(\text{Verlust}) = 1 - \left(\frac{285}{360}\right)^3 \approx 50.38 \%$$
  
}%% EN TNT
%%\newpage


\subsubsection{Gesetze}
Um Wahrscheinlichkeiten aufzuzeichnen, bietet sich das Baumdiagramm an.
Dabei zeichnet man einen «Baum» von links-nach-rechts oder wie in der Informatik üblich von oben-nach-unten.

Zu den Ästen schreibt man deren Wahrscheinlichkeiten.

Man beginnt bei der Wurzel und für jeden möglichen Ausgang zeichnet man einen Ast. Dabei gelten die folgenden Gesetze:



\begin{gesetz}{Produktregel/Pfadregel}{}\index{Produktregel!Baumdiagramm}\index{Pfadregel!Baumdiagramm}\label{experimentePfadUndSummenregel}
Alle Wahrscheinlichkeiten von der Wurzel bis zum Endknoten werden aufmultipliziert, um die Endwahrscheinlichkeit (am Endknoten) zu erhalten.
\end{gesetz}

\leserluft\leserluft
\begin{gesetz}{Summenregel I}{}\index{Summenregel!Baumdiagramm}
Alle von einem Knoten ausgehenden Äste haben in der Summe die Wahrscheinlichkeit 1 (= 100\%).
\end{gesetz}

\begin{gesetz}{Summenregel II}{}
Alle Endwahrscheinlichkeiten (bei den Endknoten\footnote{Endknoten werden auch als Blätter des Baumes bezeichnet.}) zusammen addiert, ergeben in der Summe die Wahrscheinlichkeit 1 (= 100\%).
\end{gesetz}

\begin{gesetz}{Summenregel III}{}
  Um die Wahrscheinlichkeit eines Ereignisses zu berechnen, werden die
  Wahrscheinlichkeiten all derjenigen Endknoten zusammengezählt, die zum gewünschten Ereignis gehören.
\end{gesetz}


\begin{bemerkung}{Und vs. Oder}{}
  Ein «Und» in einem Text entspricht in der Regel der Multiplikation
  während das «Oder» einer Addition entspricht:

  $$\text{ «... und ...» } \stackrel{{}^{}}{=}  \cdot{}$$
  $$\text{ «... oder ...» } \stackrel{{}^{}}{=}  +$$
\end{bemerkung}
\newpage


\subsubsection{Referenzaufgaben}

\paragraph{Urne ohne Zurücklegen} In einer Urne liegen drei grüne, zwei blaue und eine rote Kugel. Ich ziehe blind (Urne) zwei Kugeln hintereinander, ohne diese wieder zurückzulegen.

Wie groß ist die Wahrscheinlichkeit, dass genau eine blaue Kugel dabei ist?

Zeichnen Sie das zugehörige Baumdiagramm:

\TNTeop{
\bbwCenterGraphic{15cm}{geso/stoch/img/UrneBaum.png}
  
  Lösung: $P(E)=\frac8{15}=0.5333...$\vspace{13cm}}%% END TNTeop
%%\newpage

\paragraph{Ziege}\index{Ziege} Heidi und Peter spielen mit einem Würfel um eine Ziege.

Sie vereinbaren folgende Regeln:

\begin{enumerate}
\item Die Ziege erhält sofort, wer eine 5 oder eine 6 würfelt.
\item Zuerst würfelt Heidi, dann (falls Heidi noch nicht gewonnen hat) Peter, dann allenfalls noch einmal Heidi.
\item Ist nach insgesamt drei Würfen noch nichts entschieden, dann erhält Peter die Ziege.
\end{enumerate}


Zeichnen Sie das zugehörige Baumdiagramm und entscheiden Sie, ob es sich um ein faires Spiel handelt:

\TNTeop{\bbwCenterGraphic{15cm}{geso/stoch/img/HeidiPeter.png}

Peter gewinnt also mit $\frac29 + \frac8{27} = \frac{14}{27}$, was
etwas mehr als 50\% ausmacht. Das Spiel ist nicht fair, denn wer als
zweiter würfeln darf, hat eine leicht höhere Gewinnchance.}%% END TNT eop
%%\newpage
\newpage

\textbf{Aufgabe Schweizer Münzen (optional)}

In zwei Urnen sind Münzen. In jeder Urne kommt jede Schweizer Münzart
genau einmal vor:

0.05, 0.10, ... 5.--.

Aus jeder Urne wird zufällig eine Münze gezogen.

Wie groß ist die Wahrscheinlichkeit:
\begin{itemize}
\item Einen ganzen Frankenbetrag zu ziehen?
\item Mehr als 3 Franken zu erwischen?
\end{itemize}

  Tipp: Baumdiagramm oder Tabelle (Zeilen = erste Urne; Spalten =
  zweite Urne)

  \TNTeop{
    Ganzer Frankenbetrag:

    Die erste Verzweigung im Baum hat drei Äste. Die Münzen 0.05, 0.10
    und 0.20 haben je 1/7 Wahrscheinlichkeit und daraus resultiert nie
    ein ganzer Frankenbetrag. Die Münze 0.50 hat 1/7
    Wahrscheinlichkeit, aber nur ein einziger weiterführender Ast
    liefert einen ganzen Frankenbetrag, nämlich nochmals 1/7.
    Die ganzen Frankenbeträge (1.-, 2.- und 5.-) haben zusammen auch
    3/7 Wahrscheinlichkeit, aber nur die ganzen Frankenbeträge aus der
    2. Urne liefern einen ganzen Frankenbetrag.

    Somit:

    $$P(X= \text{ ganzer Frankenbetrag } ) =
    \frac17\cdot{}\frac17+\frac37\cdot{}\frac37 = \frac{22}{49}
    \approx 44.90\%$$

    Mehr als CHF 3.-:

    Wenn der 5-Liber dabei ist, so sind immer mehr als 3 Franken
    gezogen. Anders ist mehr als 3.- nur möglich mit 2.- + 2.-. Somit
    sind Somit gibt es 14 von 49 Varianten.

    $$P(X= \text{ Mehr als CHF 3.- } ) = \frac{14}{49} \approx 28.57\%$$
  }
  

\newpage

\subsection*{Aufgaben}

Mehrstufige Zufallsexperimente:


\aufgabenFarbe{
Ein Glücksrad hat 8 gleich grosse Sektoren mit den Nummern
1, 2, 1, 2, 3, 2, 1, 2. Nach Stillstand des Rades zeigt der oberhalb des Rades
fix montierte Pfeil auf einen Sektor mit Nummer.
\\
a) Das Rad wird einmal gedreht. Mit welcher Wahrscheinlichkeit erscheint 2?\\
b) Das Rad wird zweimal gedreht. Mit welcher Wahrscheinlichkeit sind beide
Ziffern gleich?\\
c) Das Rad wird fünfmal gedreht. Mit welcher Wahrscheinlichkeit ist die
Ziffernfolge 1 - 2 - 3 - 2 - 1?
}
\TNT{5.2}{
  a) $\frac{4}{8} = \frac{1}{2}$

  b) $(1;1), (2;2), (3;3): \frac{3}{8}\cdot{}\frac{3}{8} +
  \frac{4}{8}\cdot{}\frac{4}{8} + \frac{1}{8}\cdot{}\frac{1}{8} =
  \frac{13}{32} \approx 40.63\%  $
}%% end TNT

\aufgabenFarbe{Eine Klasse besteht aus 10 Schülerinnen und 14 Schülern. Durch das Los
werden 5 Personen dieser Klasse ausgewählt. Wie gross ist die
Wahrscheinlichkeit, dass\\
a) alle fünf Personen Schülerinnen sind\\
b) alle fünf Personen Schüler sind\\
d) unter den 5 Personen Schülerinnen und Schüler vorkommen?
}
\TNTeop{
a) $\frac{10}{24}\cdot{}\frac{9}{23}\cdot{}\frac{8}{22}
  \cdot{} \frac{7}{21} \cdot{} \frac{6}{20} = \frac{3}{506} \approx
  0.59\%$

b)$\frac{14}{24}\cdot{}\frac{13}{23}\cdot{}\frac{12}{22}
  \cdot{} \frac{11}{21} \cdot{} \frac{10}{20} = \frac{13}{256} \approx
  4.71\%$

c) $1 - \frac{3}{506} - \frac{13}{256} = \frac{125}{132} \approx 94.70\%$
}%% end TNT eop
%%%%%%%%%%%%%%%%%%%%%%%%%%%%%%%%%%%%%%%%%%%%%%%%%%%%%%%%%%%%%%%%%%%%%%

\aufgabenFarbe{
In einem Gefäss befinden sich 5 rote, 5 blaue und 5 gelbe Kugeln. Es werden
drei Kugeln aufs Mal gezogen. Wie gross ist die Wahrscheinlichkeit, dass\\
a) nur blaue Kugeln gezogen werden?\\
b) keine gelben Kugeln gezogen werden?
}%% end Aufgabenfarbe

\TNT{4}{
  a) $\frac{5}{15} \cdot{} \frac{4}{14} \cdot{} \frac{3}{13} =   \frac{2}{91} \approx 2.20\%$

  b) $\frac{10}{15} \cdot{} \frac{9}{14} \cdot{} \frac{8}{13} =   \frac{24}{91} \approx 26.37\%$
  
}%% end TNT


\aufgabenFarbe{
Ein Kasten enthält 2 schwarze, 2 rote und 6 grüne Kugeln. Es werden
nacheinander 3 Kugeln gezogen. Bestimmen Sie die Wahrscheinlichkeit, drei
verschiedene Farben zu erhalten, wenn\\
a) ohne Zurücklegen\\
b) mit Zurücklegen\\
gezogen wird.
}%% end Aufgabenfarbe

\TNTeop{
  a) $\frac{1}{5} = 20\%$

  b) $\frac{18}{125} = 14.4\%$
}%% end TNT eop
%%%

\aufgabenFarbe{Ein normaler Spielwürfel wird sechs Mal geworfen. Mit welcher
Wahrscheinlichkeit treten sechs verschiedene Augenzahlen auf?}%% end AufgabenFarbe

\TNT{8}{$\frac{6!}{6^6}\approx 1.54\%$}%% end TNT

\newpage

%%%%%%%%%%%%%%%%%%%%%%%%%%%%%%%%%%%%%%%%%%%%%%%%%%%%%%%%%%%%%%%%%%%%%%%%
%% Bernoulli - Ketten
\subsection{Bernoulli-Ketten}\index{Verteilung!Bernoulli}\index{Bernoulli}
(Binomialverteilung)

\TRAINER{Einstiegsvideo: ``EinfachMathe!''\\
  \texttt{https://www.youtube.com/watch?v=WC5O317JzW0}}

\subsubsection{Einstiegsaufgabe}

\paragraph{Genau zwei Sechser} Ein Würfel wird \textbf{dreimal} hintereinander geworfen.

Wie groß ist die Wahrscheinlichkeit, dass genau zwei Sechser dabei sind?


Zeichnen Sie das zugehörige Baumdiagramm:

\TNT{18}{
\bbwCenterGraphic{15cm}{geso/stoch/img/genau2sechser.png}


  Lösung mit Baum oder  Bernoulli-Formel (Binomialverteilung).

  Genau zwei Sechser: $\left(\frac16\right)^2\cdot{}\frac56 +
  \left(\frac16\right)^2\cdot{}\frac56 +
  \left(\frac16\right)^2\cdot{}\frac56 =3\cdot{}
  \left(\frac16\right)^2\cdot{}\frac56= \frac{5}{72}\approx 0.06944$%%
%%
  \vspace{10cm}%%
}%% END TNT
\newpage


\subsubsection{Bernoulli-Experiment}\index{Experiment!Bernoulli}\index{Bernoulli-Experiment}

\bbwCenterGraphic{5cm}{geso/stoch/img/BernoulliJakobI_phi.jpg}
\begin{center}\textit{Jakob Bernoulli I (1654-1705)}\end{center}

Obiges Würfelexperiment ist ein typisches
Bernoulli-Experiment\footnote{Jakob I Bernoulli, Schweizer Mathematiker,
  1654-1705}, d.\,h.:

\begin{definition}{Bernoulli-Experiment}{}
\begin{itemize}
\item Das Experiment wird $n$-mal durchgeführt\footnote{Daher wird dieser
Bernoulli-Prozess manchmal auch Bernoulli-Kette genannt. Quelle
\texttt{www.wikipedia.org} 2021-04-08}.
\item Das Einzelexperiment hat genau zwei Ausgänge: Erfolg/Misserfolg.
\item Jedes Einzelexperiment hat für die erfolgreichen Ausgänge die gleiche
      Wahrscheinlichkeit $p$ (somit ist die Wahrscheinlichkeit für den
      Misserfolg jeweils gleich $1-p$).
\item Die Einzelexperimente sind voneinander unabhängig.
\end{itemize}
\end{definition}

\begin{bemerkung}{Bernoulli-Experiment}{}
Typischerweise sind wir bei Bernoulli-Experimenten an der
Zufallsvariable $X$ interessiert, welche angibt, wie oft ein Erfolg
eingetreten ist: $X$ hat also die Werte 0, 1, 2, 3, 4, ..., $n$.
\end{bemerkung}

\textbf{Kein} Bernoulli-Experiment ist beispielsweise das Ziehen von zwei Bällen aus
einer Urne mit drei gelben und sechs orangefarbenen Bällen \textbf{ohne} diese
jeweils zurückzulegen; denn dabei verändern sich die
Wahrscheinlichkeiten der verbleibenden Bälle nach jedem Herausnehmen.
\newpage


\paragraph{Beispiele von Bernoulli-Experimenten}

\begin{itemize}
\item
In einer Urne liegen zehn Kugeln: vier blaue und sechs grüne. Ein
Einzelexperiment besteht darin, genau eine Kugel zu ziehen, die Farbe
zu notieren und die Kugel sogleich wieder zurückzulegen.
Das Gesamtexperiment besteht darin, fünf Einzelexperimente
durchzuführen. Jedesmal ist die Wahrscheinlichkeit, eine blaue Kugel
zu erwischen, gleich groß und wir sprechen von einem Bernoulli-Experiment.

\item Vierfaches Drehen am Glücksrad: Wie groß ist die
  Wahrscheinlichkeit, genau zweimal «Gewinn» zu erzielen, wenn pro
  Drehung ein Gewinn mit 28 \% Wahrscheinlichkeit auftritt?

\item Siebenfaches Würfeln mit einem Spielwürfel, bei dem wir
  interessiert sind, wie groß die Wahrscheinlichkeit ist, dass genau fünf mal ein kleinerer Wurf als eine \epsdice{3} gewürfelt wird.
  
\item
Es müssen aber nicht immer Würfel, Münzen oder Kugeln in Urnen sein:
Der Torwart Alex Calderoni\footnote{Alex Calderoni * 1976: Bekannter
  italienischer Torwart beim Fußballspiel (Quelle Wikipedia 2022).}
fängt womöglich einen Ball mit einer
Wahrscheinlichkeit von 38\%.
Wie groß ist die Wahrscheinlichkeit, dass er genau sechs von zehn Torschüssen
abfangen kann?
\end{itemize}

\TRAINER{Diese Aufgaben allenfalls auch aufs Stochastik-Aufgabenblatt.}
\newpage


\subsubsection{Formel von Bernoulli}\index{Bernoulli!Formel von}\index{Formel von Bernoulli}

\TRAINER{Einstiegs-Viedo: \texttt{https://www.youtube.com/watch?v=WC5O317JzW0}}


\begin{gesetz}{Bernoulli-Formel (Binomialverteilung)}{}
  $$P(X=k) = {{n}\choose {k}}\cdot{}p^k\cdot{}(1-p)^{n-k}$$


%%\renewcommand{\arraystretch}{2}
\begin{bbwFillInTabular}{rcl}
  $n$ &=& \TRAINER{Anzahl Durchführungen (Baumtiefe)}\\
  $p$ &=& \TRAINER{Gewinnwahrscheinlichkeit pro Durchführung}\\
  $k$ &=& \TRAINER{Anzahl geforderte Gewinne}\\
\end{bbwFillInTabular}
\end{gesetz}


Bei der sog. Bernoulli-Kette handelt es sich um ein
Bernoulli-Experiment
in einem mehrstufigen Baum ($n$ Versuche). Dabei interessiert
uns, wie groß die Wahrscheinlichkeit ist, dass von den beiden
möglichen Ausgängen (Treffer / Niete) einer davon \textbf{genau} $k$-mal auftritt.

\begin{bemerkung}{Binomialkoeffizient}{}
  Wenn ich den Baum ablaufe, gibt es immer $n$ Verzweigungen.

  An genau $k$
Stellen will ich Richtung «Treffer» abzweigen. Wie viele Varianten
habe ich, $k$ Stellen aus $n$ auszuwählen?

\vspace{4mm}

Es gibt \LoesungsRaumLen{40mm}{$n \choose k$} Möglichkeiten, $k$ Verzweigungen aus
$n$ auszuwählen.
\end{bemerkung}

\newpage
\begin{beispiel}{Sieben mal Würfeln}{}
Wie groß ist die Wahrscheinlichkeit, mit einem Würfel bei
siebenmaligem Werfen genau fünf mal eine der Zahlen \epsdice{1} oder
\epsdice{2} zu erreichen?
\end{beispiel}
Dabei legen wir fest (vgl. Beispiel oben: Siebenmaliges Würfeln):
\begin{itemize}

\item
  $n$ = Anzahl mögliche Durchführungen; hier 7

\item
  $p$ = Wahrscheinlichkeit, dass eine \epsdice{1} oder \epsdice{2}
  geworfen wird. Hier $p = \frac26$.\\
  Somit ist die Misserfolgsquote $= (1-p)$; hier $1-\frac26=\frac46$.


\item
  $k$ = Anzahl der gewünschten ``Treffer''; hier 5, denn wir wollen fünf
Mal eine \epsdice{1} oder eine \epsdice{2} erzielen.
\end{itemize}

Nun gilt

In obigem Beispiel ergibt sich:
\vspace{2mm}
$$P(X=5) = \LoesungsRaumLen{70mm}{{7\choose 5} \cdot{}
  \left(\frac26\right)^5 \cdot{} \left(1-\frac26\right)^{7-5}} =
\LoesungsRaumLen{60mm}{{7\choose 5} \cdot{} \left(\frac26\right)^5
  \cdot{} \left(\frac46\right)^2}$$


\textbf{Taschenrechner}:

\leserluft

Dies kann mit dem Taschenrechner gelöst werden:
Entweder direkt ...

$$\LoesungsRaumLang{( 7 \text{\,\,nCr\,\,} 5 ) *\left(\frac26\right)^5 *
\left(1-\frac26\right)^{7-5}} \approx \LoesungsRaum{0.03841 = 3.841\% =\frac{28}{729}}$$

... oder aber mit der eingespeicherten Formel bei
\tiprobutton{data_stat-reg-distr} und dort unter \texttt{DISTR: 4
  Binomialpdf: SINGLE}

... dann ...

\begin{tabular}{c@{ $:$ }l}
  $n=7$   & ... bei \texttt{n} \\
  $p=2/6$ & bei \texttt{p(\text{success})} und \\
  $k=5$   & bei \texttt{x} eintragen.
\end{tabular} 
\newpage


\subsection*{Aufgaben}

\aufgabenFarbe{Man würfelt 4-mal hintereinander mit einem gewöhnlichen Spielwürfel. Wie
gross ist die Wahrscheinlichkeit, dass man genau 2 Sechser gewürfelt
hat?}
\TNT{4}{
  $$P(X=2) = {4 \choose 2} \cdot{} \left(\frac{1}{6}\right)^2 \cdot{}
  \left(\frac{5}{6}\right)^{4-2} \approx 11.57\%$$
}%% end TNT


\aufgabenFarbe{Maturaprüfungen:\\2016 (GESO) Aufgabe 11 (Geburten)\\2017 (GESO) Aufgabe 17. (borkenkäferresistent)\\2018 (GESO) Serie 1 Aufg. 16. (Papayasendung)\\2018 (GESO) Serie 2 Aufg. 15. (Sprössling)\\2018 (GESO) Serie 3 Aufg. 14. (Ikosaeder)}


\newpage
\subsubsection{Kumulierte Wahrscheinlichkeit}

\subsubsection*{Motivation}

\begin{beispiel}{Hygienemasken}{}
  0.2\,\% der hergestellten Hygienemasken sind defekt. Ich kaufe 1000
  Stk. Wie groß ist die Wahrscheinlichkeit, dass weniger als 6 defekt sind?
  \end{beispiel}
\newpage

Basketballspieler «Basil Ballisti» trifft mit einer
Wahrscheinlichkeit von 85\%.

Die Wahrscheinlichkeit, dass er von 20 Würfen \textbf{genau} 17 Mal trifft, kennen wir
bereits von der Formel von Bernoulli (Bernoulli-Experiment/Binomialverteilung):

$$P(X=17) = \LoesungsRaumLang{{20 \choose 17}\cdot 0.85^{17}\cdot(1-0.85)^{20-17}}=$$
$$\LoesungsRaumLang{{20 \choose 17}\cdot 0.85^{17}\cdot (0.15)^{3}}\approx{}\LoesungsRaum{24.28\%}$$
Dabei bezeichnet $E$ das Ereignis ``siebzehn Treffer in den 20 Würfen''.

\begin{beispiel}{Kumulierte Wahrscheinlichkeit}{}
Wie groß ist nun aber die Wahrscheinlichkeit, dass er bei 20 Würfen
\textbf{höchstens} (= maximal)  17 mal trifft?
\end{beispiel}


Lösung: Wir summieren alle Wahrscheinlichkeiten auf, bei denen er
während seinen 20 Würfen
keinen Treffer, einen Treffer, zwei Treffer, ... bis zu 17 Treffer
erzielt. Dies nennen wir die kumulierte\footnote{«Kumulieren» kommt von
  lat. \texttt{cumulus} = Anhäufung.} Wahrscheinlichkeit\index{Wahrscheinlichkeit!kumulierte}:

$$P(X\le 17) = $$
\TNT{6}{$$P(X=0) + P(X=1) + P(X=2) + ... + P(X=16) + P(X=17)$$
$$={20\choose 0}\cdot 0.85^0 \cdot 0.15^{20} + {20\choose 1}\cdot 0.85^1 \cdot 0.15^{19} + {20\choose 2}\cdot 0.85^2 \cdot 0.15^{18} + ... $$

$$... + {20 \choose 17}\cdot{} 0.85^{17} \cdot{0.15}^3$$
  
}%% END TNT

Dies können wir auch mit dem Summenzeichen ($\Sigma$) schreiben:
$$P(X\le 17) = \LoesungsRaumLang{\sum_{k=0}^{17} {20 \choose k} \cdot 0.85^k \cdot 0.15^{20-k}}$$

\newpage
\textbf{Taschenrechner}:

\leserluft

Für den Taschenrechner verwenden Sie folgende Notation: 

$$\LoesungsRaumLang{\sum_{x=0}^{17}\left(( 20 \text{\,\,nCr\,\,} x ) * 0.85^x *
0.15^{20-x}\right)} \approx \LoesungsRaum{0.5951}$$




\begin{gesetz}{Kumulierte Wahrscheinlichkeit}{}
  $$P(X \le k_{\text{max}}) = \sum_{k=0}^{k_{\text{max}}} {n \choose k} \cdot{} p^k \cdot{} (1-p)^{n-k}$$


\begin{tabular}{rcl}
  $n$ &=& Anzahl Durchführungen (Baumtiefe)\\
  $p$ &=& Trefferwahrscheinlichkeit pro Durchführung\\
  $k_{\text{max}}$ &=& \textbf{Maximal} erreichte Treffer\\
\end{tabular}
\end{gesetz}


Doch auch hierzu gibt es eine eingespeicherte Formel
bei
\tiprobutton{data_stat-reg-distr} und dort unter \texttt{DISTR: 5
  Binomial\textbf{\color{red} c}df: SINGLE}

dann

$n=20$ bei $n$,

$p=0.85$ bei $p(\text{success})$ und

$k_{\text{max}}=17$ bei $x$ eintragen.

\newpage

\subsection*{Aufgaben}

\aufgabenFarbe{Wir wissen: 0.2\,\% der hergestellten Hygienemasken
  sind defekt. Wie groß ist die Wahrscheinlichkeit, dass von zufällig
  ausgewählten 1000 Stück maximal 5 defekt sind?}

\TNT{2}{$p=0.002$, $n=1000$, $P(X \le 5) \approx 98.35\,\%$}

\aufgabenFarbe{Wie groß ist die Wahrscheinlichkeit, dass
  Cristiano Ronaldo (*1985, Fußballspieler), dessen durchschnittliche Trefferwahrscheinlichkeit bei 84\% liegt,
  in einem Training mit 40 Torschüssen weniger als 36 Tore erzielt?
}
\TNT{2}{
$p=0.84$; $n=40$; $k=35$ somit $P(X < 36) = P(X\le 35) \approx 78.90\,\%$
}


\aufgabenFarbe{In einer Stadt sind 4.5\% der Autos rot. Sie stehen an der Ampel und betrachten 25 vorbeifahrende Autos.\\
  a) mit welcher Wahrscheinlichkeit ist kein rotes Auto dabei?\\
  b) mit welcher Wahrscheinlichkeit ist genau ein rotes Auto dabei?\\
  c) mit welcher Wahrscheinlichkeit sind maximal 3 rote Autos dabei (also weniger als 4 rote Autos)?
}

\TNT{4}{
    a) $31.63\%$ (binomialpdf mit $n=25; p=0.045; x=0$)\\
    b) $37.25\%$ (binomialpdf mit $n=25; p=0.045; x=1$)\\
    c) $97.57\%$ (binomial\textbf{c}df mit $n=25; p=0.045; x=3$)\\
}%% end TNT
\newpage


\aufgabenFarbe{Die Herstellung von LEDs (Licht emittierende Dioden) hat einen Ausschuss (= defekte Teile) von 0.5\%. Sie kaufen 30 solcher Dioden (LEDs) für ein Nachtlicht.
  Mit welcher Wahrscheinlichkeit sind maximal zwei davon defekt?}
\TNT{3.2}{
  $$P(X \le 2) = \sum_{k=0}^{2}{30 \choose k} \cdot{}0.005^k \cdot{} 0.995^{30-k} \approx 99.954\%$$
}%% end TNT

%%%%%%%%%%%%%%%%%%%%%%%%%%%%%%%%%%%%%%%%%%%%%%%%%%


\aufgabenFarbe{
Wie groß ist die Wahrscheinlichkeit mit einem Spielwürfel
(\epsdice{1} bis \epsdice{6}) in 10 Würfen...\\
a) ... genau viermal eine \epsdice{6} zu werfen?\\
b) ... maximal eine \epsdice{6} zu werfen?\\
c) ... nicht viermal eine eine \epsdice{6} zu werfen?\\
d) ... mehr als dreimal eine \epsdice{6} zu werfen?\\
}%% end Aufgabenfarbe
\TNTeop{
a) $P(X=4) = {10 \choose 4} \cdot{} \left( \frac16 \right)^4 \cdot{}
\left(\frac56\right)^6 \approx 5.43\%$

b) $P(X \le 1) =\sum_{k=0}^{1} {10 \choose k} \cdot{} \left( \frac16 \right)^k \cdot{}
\left(\frac56\right)^{10-k} \approx 48.45\%$

c) $P(X \ne 4) = 1- P(X = 4) \approx 1 - 0.0543 = 94.57\%$

d) $P(X > 3) = 1-P(X\le 3) = 1 - \sum_{k=0}^{3} {10 \choose k} \cdot{} \left( \frac16 \right)^k \cdot{}
\left(\frac56\right)^{10-k} \approx 6.97\%$

}%% end TNTeop (implicit \newpage)


%% Gegenwahrscheinlichkeit:
\aufgabenFarbe{Ein Frühstücksflocken-Hersteller behauptet, in jeder 10. Packung ist ein Spielzeug enthalten. Sie kaufen 10 Packungen.\\
  a) Mit welcher Wahrscheinlichkeit haben Sie mindestens zwei Spielzeuge in Ihren 10 Packungen?\\
  b) (*) Wie viele Packungen müssen Sie kaufen, damit Sie mit über 95\% Wahrscheinlichkeit mindestens zwei Spielzeuge dabei haben?}

\TNT{4}{
  Gegenwahrscheinlichkeit:
  
  a) $$P(X \ge 2) = 1 - P(X \le 1) = 1 - 73.61\% \approx 26.39\%$$
  b) 46 Stück. Pröbeln Sie mit binomialcdf mit
  $p=0.1$ und  $x=1$ solange mit $n$, bis die Wahrscheinlichkeit maximal 1 Spielzeug zu erreichen kleiner als 5\% wird. Dies ist erstmals bei 46. Packungen der Fall.
}%% end TNT



\aufgabenFarbe{Lukas bringt zu einem Würfelspiel seinen eigenen Spielwürfel mit. Im Verlauf
des Spiels entsteht der Verdacht, Lukas’ Würfel zeige überdurchschnittlich
viele Sechsen an.\\
Man beschliesst, einen Versuch durchzuführen: Es wird 30-mal gewürfelt. Es
wird vereinbart: Wenn 10-mal oder öfter eine Sechs erscheint, erklären wir
Lukas’ Würfel als gefälscht.\\
Lukas sagt: «Selbst wenn der Test zehn oder mehr Sechsen ergibt, heisst das
keineswegs, dass mein Würfel gefälscht ist; das Resultat kann auch bei einem
unverfälschten Würfel eintreten.»\\
Wie gross ist diese Wahrscheinlichkeit, dass
in 30 Würfen ein gerechter Würfel 10 oder mehr Sechsen anzeigt?
}

\TNTeop{
 $P(X \ge 10) = \sum_{k=10}^{30} {30 \choose k} \cdot{} \left( \frac16 \right)^k \cdot{}
  \left(\frac56\right)^{30-k} \approx 1.97\%$

  Mit ca. $2\%$ Wahrscheinlichkeit tritt auch bei einem gerechten Würfel ein derart
extremes Ergebnis auf. Ein gerechter Würfel wird mit ca. 98\% Wahrscheinlichkeit
ein moderateres Ergebnis zeigen (weniger als 10 Sechsen). Man wird also
eher an eine Fälschung glauben. Wer sagt, es sei ein gefälschter
Würfel irrt sich nur mit ca. $2\%$ Wahrscheinlichkeit (Hier:
\textbf{Irrtumswahrscheinlichkeit = Fehlerwahrscheinlichkeit} $\approx 2\%$).
}%% end TNTeop


\newpage

%%%%%%%%%%%%%%%%%%%%%%%%%%%%%%%%%%%%%%%%%%%%%%%%%%%%%%%%%%%%%%%%%%%%%%%%

\subsection{Lotto-Wahrscheinlichkeit}\index{Verteilung!Hypergeometrische}\index{Hypergeometrische Verteilung}\index{Lotto-Wahrscheinlichkeit}
In einer Urne liegen sieben Kugeln. Drei davon sind Treffer (grün) und vier davon sind Nieten (schwarz).
Wir dürfen zwei Kugeln blind herausnehmen (ohne diese wieder zurückzulegen).

\bbwCenterGraphic{5cm}{geso/stoch/img/Urne3T4N.png}

Wie groß ist nun die Wahrscheinlichkeit
\begin{itemize}
\item beide Treffer zu erzielen ($t=2$),
\item genau einen Treffer zu erzielen ($t=1$) oder
\item gar keinen Treffer zu erzielen ($t=0$)?
\end{itemize}

Diese Wahrscheinlichkeiten können wir nun einerseits durch ein Baumdiagramm lösen, indem wir die Treffer mit \textbf{\color{ForestGreen}T} und die Nieten mit \textbf{\color{red}N} bezeichnen:

\noTRAINER{\mmPapier{7.2}}
\TRAINER{
\bbwCenterGraphic{7cm}{geso/stoch/img/Baum2aus3T4N.png}
}

Somit erhalten wir für zwei Treffer ($t=2$) eine Wahrscheinlichkeit von $\frac{6}{42}$; für einen Treffer ($t=1$) die Summe aus $\frac{12}{42}$ und nochmals $\frac{12}{42}$ also total $\frac{24}{42}$ und schlussendlich für keinen Treffer ($t=0$) die Wahrscheinlichkeit $\frac{12}{42}$.
Die Summe aller Wahrscheinlichkeiten muss dann immer gleich 1.0 (100\%) sein.
\newpage

Andererseits gibt es dazu auch eine Formel, die
\textbf{Lotto-Wahrscheinlichkeit} oder «\textit{hypergeometrische
  Verteilung}»\footnote{Die Folgeglieder der \textit{geometrischen
    Verteilung} $P(X=t) = p\cdot{}q^{t-1}$ verhalten sich wie eine
  geometrische Reihe (die Quotienten von zwei Folgengliedern $P(X=t) :
  P(X=t+1)$ sind jeweils konstant). Da die hier betrachtete Verteilung
  noch rascher ``abfällt'', wie die \textit{geometrische Verteilung},
  so sprechen wir hier von \textbf{hyper}geometrisch.}.

\begin{beispiel}{Lotto-Wahrscheinlichkeit}{}
  In obiger Urne mit \textbf{drei} Treffern und \textbf{vier} Nieten
  ziehen wir genau \textbf{zwei} Kugeln heraus. Wie groß ist die
  Wahrscheinlichkeit auf zwei verschiedene Kugeln?
\end{beispiel}



Es seien:
\begin{itemize}
\item $T$ die Anzahl der Treffer in der Urne. Hier die Anzahl der grünen Kugeln (also $T = 3$)
\item $N$ die Anzahl der Nieten in der Urne. Hier die Anzahl schwarze Kugeln ($N=4$)
\item $T+N$ die Anzahl Elemente (Treffer + Nieten) in der Urne (hier $T+N = 7$)

\item $t$ die Anzahl der gewünschten bzw. zu erreichenden Treffer. Hier \zB $t=1$.
  
\item $n$ die Anzahl «gewünschten» Nieten (hier z. B. $n = 1$)
\item $t+n$ die Anzahl total herausgezogener Kugeln (hier $t+n=1+1=2$)
\end{itemize}

\newpage


\begin{gesetz}{}{}
$$P(X=t) = \frac{ {T\choose t} \cdot{} {{N}\choose {n}}}{{{T+N} \choose {t+n}}}$$
\end{gesetz}

Wie groß ist die Wahrscheinlichkeit, beim Zahlenlotto (Euromillions)
genau 3 Treffer aus 50 zu erzielen, wenn 5 aus 50 gezogen worden sind.

$$P(X=3) = \LoesungsRaumLen{100mm}{\frac{ {5 \choose 3} \cdot{} {45
      \choose 2} }{{50 \choose 5}} \approx 0.467 \%}$$

\newpage

\subsection*{Aufgaben}

\aufgabenFarbe{
Robin kauft im Supermarkt 17 Artikel ein und scannt den Einkauf
selbständig. Robins fünfjähriges Kind «schmuggelt» nachträglich zwei
Tafeln Schokolade in die Einkaufstasche. Somit sind nun 19 Artikel in
der Einkaufstasche, wovon zwei nicht abgescannt wurden.
\\
Bei einer Kontrolle durch das Ladenpersonal werden jeweils fünf
Artikel wahllos aus den Einkaufstaschen herausgepickt und mit der
Abrechnung verglichen.
\\
a) Wie groß ist die Wahrscheinlichkeit, dass bei der Prüfung genau eine der nicht
abgescannten Schokoladentafeln dabei ist?
\\
b) (optional) Wie groß ist die Wahrscheinlichkeit, dass bei der Prüfung
mindestens eine der nicht abgescannten Schokoladentafeln dabei ist?
}%% end AufgabenFarbe

\TNT{8}{
  a) Lotto Wahrscheinlichkeit. $T=2; N=17; t=1; n=4$
  $$
        P(X=1) = \frac{{2 \choose 1} \cdot{} {17 \choose 4}}{{19 \choose 5}} =
\frac{70}{171}\approx  40.95\%$$


b) $$P(x > 1) = P(x=1) + P(x=2) = \frac{70}{171} +
    \frac{ {2 \choose 2} \cdot{} {17 \choose 3} }{{19 \choose 5}} = \frac{70}{171} + \frac{10}{171}
    = \frac{80}{171} \approx 46.78\% $$
}%% end TNT



\olatLinkGESOKompendium{5.5.3}{49}{25. (Herz-Karten), 26. (Smarties), 24. (Zahlenlotto) und optional
  27. (Euromillions)}
\TRAINER{Auch noch weitere Aufgaben sammeln (S. Sammlung M. Bräm).}
\newpage
